\documentclass[aspectratio=169,8pt]{beamer}
\usetheme{Madrid}
\usepackage[utf8]{inputenc}
\usepackage[T1]{fontenc}
\usepackage{amsmath,amssymb}
\usepackage{booktabs}
\usepackage{enumitem}
\setlist[enumerate]{label=\Alph*.,leftmargin=*,itemsep=1pt,topsep=2pt}
\setlist[itemize]{leftmargin=*,itemsep=1pt,topsep=2pt}
\setbeamerfont{block body}{size=\tiny}
\setbeamerfont{block title}{size=\scriptsize}
\setbeamerfont{frametitle}{size=\footnotesize}
\setbeamerfont{normal text}{size=\tiny}
\setbeamertemplate{navigation symbols}{}
\setbeamertemplate{itemize item}{\tiny$\bullet$}
\setbeamertemplate{itemize subitem}{\tiny$\circ$}

\title[GARP RAI]{GARP Risk \& AI --- 250 Detailed Hard Questions}
\subtitle{Unofficial Exam Prep}
\author{Unofficial}
\date{\today}

\begin{document}
	\begin{frame}\titlepage\end{frame}
	
	% =================================================================
	\section{Module 1: AI and Risk --- Introduction \& Overview}
	% =================================================================
	
	\begin{frame}{Q1 --- Inscrutability vs Explainability in Practice}
		\begin{block}{Question}
			A tier-1 bank has deployed a deep neural network to screen consumer loan applications. When a rejected applicant files a formal complaint under GDPR's ``right to explanation'', the compliance team asks the data science team for a justification. The data science team can run a SHAP explainer that produces the top three features driving the specific rejection, but when asked to describe how the model internally transforms the input vector into the output probability, they admit that not even they can trace this path. The regulator accepts the SHAP output as a valid post-hoc explanation. Which statement most accurately characterises this scenario?
		\end{block}
		\begin{enumerate}
			\item The model is interpretable but not explainable, because SHAP values are a property of the model rather than the specific decision.
			\item The model is explainable but not interpretable, because post-hoc explanation of a specific decision is possible even though the internal mechanism cannot be described.
			\item The model is neither explainable nor interpretable, because SHAP is a post-hoc technique that does not reveal the model's true behaviour.
			\item The model is both fully interpretable and explainable, because SHAP output is equivalent to understanding the model's internal mechanism.
		\end{enumerate}
		\begin{exampleblock}{Answer: B}\end{exampleblock}
		\begin{block}{Explanation}
			\textbf{Explainability} refers to the ability to provide a human-understandable account of a specific model decision --- typically post-hoc, using techniques like SHAP or LIME. \textbf{Interpretability} refers to the ability to understand the model's internal logic and mechanisms. Deep neural networks are frequently explainable but not interpretable. Option A is wrong because SHAP does not make the model internally interpretable. Option C is wrong because post-hoc explanation \emph{does} count as explainability. Option D is wrong because SHAP values do not reveal the model's internal mechanism; they only attribute the output to inputs.
		\end{block}
	\end{frame}
	
	\begin{frame}{Q2 --- Inscrutability in AI Models}
		\begin{block}{Question}
			The risk management committee at a global insurance group is reviewing the group's AI inventory. The CRO notes that one class of models in the fraud-detection stack is described in the documentation as ``inscrutable''. A junior analyst asks what this means in the context of the GARP Risk and AI curriculum. The senior data scientist explains that inscrutability refers to a specific limitation of certain AI models. Which of the following most accurately captures the meaning of inscrutability as used in the GARP curriculum?
		\end{block}
		\begin{enumerate}
			\item Inscrutability refers to the inability of the model to generalise to out-of-sample data.
			\item Inscrutability refers to the inability to understand the internal mechanisms through which the model produces its outputs.
			\item Inscrutability refers to the inability to explain a specific prediction to a customer or regulator.
			\item Inscrutability refers to the inability to measure the model's accuracy with traditional statistical metrics.
		\end{enumerate}
		\begin{exampleblock}{Answer: B}\end{exampleblock}
		\begin{block}{Explanation}
			Inscrutability is specifically the property of a model whose internal mechanisms are opaque --- the analyst cannot trace how the input vector is transformed into the output. Option A confuses inscrutability with generalisability. Option C confuses it with explainability, which is about specific decisions. Option D confuses it with accuracy or measurability.
		\end{block}
	\end{frame}
	
	\begin{frame}{Q3 --- Classical AI versus Deep Learning in Practice}
		\begin{block}{Question}
			An asset manager is considering two ways to build a market-regime classifier. Approach One is a rule-based system that uses a fixed set of hand-coded conditions on moving averages and volatility thresholds to label regimes as ``risk-on'' or ``risk-off''. Approach Two trains a deep neural network with several hidden layers on historical features to output the same classification. The CIO asks you which approach corresponds to which paradigm in AI and what the key trade-offs are. Which statement best captures the distinction?
		\end{block}
		\begin{enumerate}
			\item Approach One is deep learning and Approach Two is classical AI, because deep neural networks rely on pre-defined rules while rule-based systems learn from data.
			\item Approach One is classical AI and Approach Two is deep learning; the classical approach scales well to complex unstructured data, while the deep learning approach does not.
			\item Approach One is classical AI (rules and search within a well-defined context) and Approach Two is deep learning (learns representations via backpropagation); the deep learning approach can capture non-linear patterns but is typically more inscrutable.
			\item Both approaches are classical AI because both encode a mapping from inputs to outputs that is fixed after deployment.
		\end{enumerate}
		\begin{exampleblock}{Answer: C}\end{exampleblock}
		\begin{block}{Explanation}
			Classical AI (GOFAI) uses hand-crafted rules and symbolic reasoning. Deep learning learns from data using backpropagation and gradient descent, allowing it to capture non-linear relationships, but the internal mechanism is opaque (inscrutable). Option A reverses the definitions. Option B also reverses the scalability argument. Option D ignores the fact that deep learning learns its parameters from data, which is definitionally distinct from hand-coded rules.
		\end{block}
	\end{frame}
	
	\begin{frame}{Q4 --- Deep Learning's Breakthroughs}
		\begin{block}{Question}
			A history-of-AI seminar traces three milestones: (i) Deep Blue's 1997 defeat of Kasparov, (ii) AlexNet's 2012 victory in the ImageNet competition, and (iii) the 2017 publication of ``Attention Is All You Need''. A candidate is asked to associate each milestone with the correct AI paradigm. Which of the following correctly maps the milestones?
		\end{block}
		\begin{enumerate}
			\item (i) deep learning, (ii) classical AI, (iii) classical AI.
			\item (i) classical AI, (ii) deep learning, (iii) foundational work for modern LLMs via the transformer architecture.
			\item (i) classical AI, (ii) reinforcement learning, (iii) rule-based expert systems.
			\item (i) deep learning, (ii) transformer-based models, (iii) reinforcement learning.
		\end{enumerate}
		\begin{exampleblock}{Answer: B}\end{exampleblock}
		\begin{block}{Explanation}
			Deep Blue was a classical AI system based on search and hand-crafted evaluation. AlexNet is a landmark deep convolutional neural network that won ImageNet in 2012, launching the modern deep learning era. ``Attention Is All You Need'' introduced the transformer, which underpins modern large language models. Option A and D misattribute. Option C mislabels AlexNet and the transformer paper.
		\end{block}
	\end{frame}
	
	
	\begin{frame}{Q6 --- Deep Learning and Unstructured Data}
		\begin{block}{Question}
			A risk analyst observes that a hedge fund's new AI-driven news-parsing engine routinely outperforms a legacy keyword-based NLP pipeline on the task of extracting sentiment from unstructured earnings-call transcripts. The legacy system relies on a hand-curated sentiment lexicon; the new system uses a pretrained transformer fine-tuned on the firm's transcripts. Which of the following best explains the performance gap in terms of the classical-AI versus deep-learning distinction?
		\end{block}
		\begin{enumerate}
			\item The legacy system fails because it cannot process numerical data, whereas the transformer can handle both text and numbers.
			\item The legacy system cannot adapt to the context in which words appear; the transformer's attention mechanism captures long-range dependencies and polysemy, allowing it to learn patterns that hand-crafted lexicons miss.
			\item The legacy system is more accurate in theory but is hampered by bugs; the transformer has no such limitation.
			\item The performance gap is due solely to the larger dataset available to the transformer, not to any architectural difference.
		\end{enumerate}
		\begin{exampleblock}{Answer: B}\end{exampleblock}
		\begin{block}{Explanation}
			Classical NLP pipelines based on lexicons fail on context-dependent meaning, negation, sarcasm, and multi-word expressions. Transformer-based models use self-attention to build contextual embeddings that capture long-range dependencies and polysemy, allowing them to learn patterns that hand-crafted dictionaries miss. Option A is factually wrong. Option C is not a plausible explanation. Option D is incomplete: data matters, but the architectural difference is decisive.
		\end{block}
	\end{frame}
	
	\begin{frame}{Q7 --- Data Types in Risk Data}
		\begin{block}{Question}
			A quantitative risk analyst is preparing a feature set for a credit-scoring model. The dataset contains the following five fields: (i) an internal credit rating on a scale of 1--10 from lowest to highest risk; (ii) the applicant's industry code from a list of 20 mutually exclusive industry classifications; (iii) the applicant's annual revenue in USD; (iv) the applicant's age in years; (v) the applicant's self-reported ESG score on a scale of 1--5 where 1 is ``very poor'' and 5 is ``very good'' but the intervals are not known to be equal. Which pair correctly classifies (i) and (ii)?
		\end{block}
		\begin{enumerate}
			\item (i) is nominal; (ii) is ordinal.
			\item (i) is ordinal; (ii) is nominal.
			\item (i) is interval; (ii) is ratio.
			\item (i) is nominal; (ii) is interval.
		\end{enumerate}
		\begin{exampleblock}{Answer: B}\end{exampleblock}
		\begin{block}{Explanation}
			(i) is an internal credit rating on an ordered scale with unequal intervals between levels --- this is ordinal. (ii) is an industry code with no inherent ordering between categories --- this is nominal. Options A, C, and D misclassify one or both fields.
		\end{block}
	\end{frame}
	
	\begin{frame}{Q8 --- Standardization vs Normalization}
		\begin{block}{Question}
			A modelling team is preparing to fit a KNN-based fraud detector on features that include ``transaction amount in USD'' (range 1 to 1,000,000) and ``days since account opened'' (range 1 to 3,650). A junior analyst proposes min-max normalization on both features. A senior analyst disagrees and proposes z-score standardization. Which of the following statements best captures the trade-off?
		\end{block}
		\begin{enumerate}
			\item Min-max normalization should always be preferred because it produces bounded inputs that stabilize distance computations.
			\item Standardization should be preferred when the feature distribution contains extreme outliers, because min-max normalization compresses the bulk of the distribution into a narrow sub-interval of [0,1].
			\item Standardization should always be preferred regardless of the shape of the distribution, because it is distribution-free.
			\item Neither scaling method should be applied, because KNN is invariant to monotone transformations of the features.
		\end{enumerate}
		\begin{exampleblock}{Answer: B}\end{exampleblock}
		\begin{block}{Explanation}
			Min-max normalization is sensitive to outliers because the min and max are affected by extreme values; the rest of the distribution is squeezed into a narrow band. Standardization uses mean and standard deviation, which are less sensitive to outliers, and preserves the shape of the underlying distribution. Option A is wrong because min-max is exactly the wrong choice under heavy tails. Option C is wrong because standardization is not distribution-free in effect; it is still a linear transformation. Option D is wrong because KNN is not invariant to monotone transformations; distances change.
		\end{block}
	\end{frame}
	
	\begin{frame}{Q9 --- Stepwise Regression and AIC}
		\begin{block}{Question}
			An analyst runs forward stepwise regression on a dataset of annual salaries using polynomial terms in age. At step two, the model ``Age + Age$^3$'' achieves AIC = 3123.20. At step three, the model ``Age + Age$^3$ + Age$^5$'' achieves AIC = 3124.50. The analyst must report which model to select and why. Which of the following is the correct conclusion?
		\end{block}
		\begin{enumerate}
			\item Select ``Age + Age$^3$ + Age$^5$'' because adding more terms cannot worsen the model's fit.
			\item Select ``Age + Age$^3$'' because AIC penalises complexity, and the increase from 3123.20 to 3124.50 indicates that the additional term is not justified.
			\item Select ``Age + Age$^3$ + Age$^5$'' because AIC is monotone in the number of parameters and only the direction of change matters.
			\item Neither model is valid because AIC cannot be compared across models with different numbers of parameters.
		\end{enumerate}
		\begin{exampleblock}{Answer: B}\end{exampleblock}
		\begin{block}{Explanation}
			AIC explicitly penalises additional parameters. In forward stepwise regression, we add a term only if it reduces AIC. Since adding Age$^5$ increased AIC from 3123.20 to 3124.50, the term is not justified. Option A is wrong because AIC, unlike R$^2$, can worsen when unnecessary terms are added. Option C confuses AIC with R$^2$. Option D is wrong because AIC is designed for cross-model comparison.
		\end{block}
	\end{frame}
	
	\begin{frame}{Q10 --- ML vs Classical Statistics}
		\begin{block}{Question}
			A research team is asked to test the hypothesis that the CAPM beta of a particular equity is significantly different from one. Separately, another team is asked to build a model that maximises the out-of-sample prediction accuracy of next-day returns for the same equity. Which of the following statements correctly matches each task to the appropriate methodology?
		\end{block}
		\begin{enumerate}
			\item Both tasks should use machine learning, because ML is always superior when the data is large.
			\item Both tasks should use classical statistics, because ML cannot handle financial time-series data.
			\item The first task is best served by classical statistics (inference, t-test on $\beta$); the second task is best served by machine learning (regularised regression, tree ensembles, or neural networks optimised for out-of-sample prediction).
			\item The first task is best served by machine learning; the second by classical statistics.
		\end{enumerate}
		\begin{exampleblock}{Answer: C}\end{exampleblock}
		\begin{block}{Explanation}
			The first task is a hypothesis test --- this is inference, the domain of classical statistics. The second task is a prediction problem --- this is the domain of machine learning. Option A is wrong because ML's strengths lie in prediction, not inference. Option B is wrong because ML handles time-series data well when used appropriately. Option D reverses the correct mapping.
		\end{block}
	\end{frame}
	
	\begin{frame}{Q11 --- Feature Scaling in Distance-Based Models}
		\begin{block}{Question}
			A junior quant trains a KNN classifier on a dataset with two features: ``customer tenure in months'' (values between 1 and 240) and ``monthly income in USD'' (values between 800 and 200,000). Without scaling, the classifier's predictions are almost entirely driven by income. She then applies z-score standardization to both features and the model performance improves dramatically. Which of the following best explains what happened?
		\end{block}
		\begin{enumerate}
			\item The KNN algorithm was broken before scaling; standardization fixed a bug.
			\item Because KNN uses a distance metric, features with larger numerical ranges dominate the distance computation; standardization brings both features onto a comparable scale so that tenure can contribute appropriately.
			\item Standardization increased the amount of information in the dataset.
			\item Standardization eliminated outliers that were causing the poor performance.
		\end{enumerate}
		\begin{exampleblock}{Answer: B}\end{exampleblock}
		\begin{block}{Explanation}
			Distance-based algorithms (KNN, K-means, SVM, PCA) are sensitive to feature scale. A feature measured in the tens of thousands (income) dominates one measured in the tens or hundreds (tenure). Standardization puts both features on the same scale (mean zero, SD one), allowing tenure to contribute. Option A is wrong --- the algorithm was not broken. Option C is wrong --- scaling does not add information. Option D is wrong --- standardization does not remove outliers; it merely scales.
		\end{block}
	\end{frame}
	
	
	\begin{frame}{Q13 --- Log Transformation for Skewed Data}
		\begin{block}{Question}
			A risk analyst examines the distribution of a feature representing ``individual trade size in USD''. The minimum is \$100 and the maximum is \$14 million. The ratio of max to min is approximately 140,000. A histogram shows a very long right tail. The analyst decides to log-transform the feature. Which of the following best describes the effect of this decision?
		\end{block}
		\begin{enumerate}
			\item The log transformation eliminates outliers from the dataset.
			\item The log transformation typically reduces positive skewness and stabilises the variance of the feature, making it more suitable as an input to linear models that assume well-behaved residuals.
			\item The log transformation converts the feature into a categorical variable.
			\item The log transformation has no effect on the shape of the distribution; only scaling methods change the shape.
		\end{enumerate}
		\begin{exampleblock}{Answer: B}\end{exampleblock}
		\begin{block}{Explanation}
			Log transformations reduce right skewness and stabilise variance in features with multiplicative structure (common in financial data like trade sizes, incomes, or market caps). Option A is wrong --- logs do not delete data points. Option C is wrong --- logs keep the feature numeric. Option D is wrong because log transformations do change the shape, unlike scaling.
		\end{block}
	\end{frame}
	

	
	\begin{frame}{Q15 --- Outliers and Model Sensitivity}
		\begin{block}{Question}
			A risk team is comparing three candidate models to predict the severity of operational losses: (i) ordinary least squares regression with a linear specification, (ii) a decision tree regressor, and (iii) a support vector machine with an RBF kernel. The dataset contains several extreme losses caused by one-off events. Which statement best captures the relative robustness of these three models to outliers?
		\end{block}
		\begin{enumerate}
			\item All three models are equally sensitive to outliers because they all minimize squared error.
			\item OLS is the most sensitive because squared residuals heavily penalise large deviations; the tree and the SVM are typically more robust because they do not rely on a globally quadratic loss.
			\item The tree is the most sensitive because it splits on thresholds, which are distorted by extreme values.
			\item The SVM is the most sensitive because the RBF kernel amplifies distances.
		\end{enumerate}
		\begin{exampleblock}{Answer: B}\end{exampleblock}
		\begin{block}{Explanation}
			OLS uses a squared loss that heavily penalises large residuals, so a single extreme observation can substantially distort the fitted line. Decision trees split on thresholds and are less sensitive to the exact magnitude of a single point. SVMs with margins are also relatively robust. Option A is wrong because the loss functions differ. Option C is wrong because trees are typically more robust to outliers, not less. Option D is wrong because the RBF kernel is a modelling choice; the SVM's robustness comes from its margin formulation.
		\end{block}
	\end{frame}
	
	\begin{frame}{Q16 --- K-Fold Cross-Validation Setup}
		\begin{block}{Question}
			A team building a credit-default classifier has 1,000 labelled observations with a 95/5 class imbalance. They consider three evaluation schemes: (i) a single 70/15/15 train/validation/test split, (ii) 10-fold cross-validation, (iii) 10-fold \emph{stratified} cross-validation. Which scheme is most appropriate, and why?
		\end{block}
		\begin{enumerate}
			\item Scheme (i), because a single split avoids leakage.
			\item Scheme (ii), because k-fold uses all data.
			\item Scheme (iii), because stratification ensures each fold preserves the 95/5 class ratio, avoiding folds with zero minority-class examples and yielding a less biased performance estimate.
			\item None of the schemes are appropriate; oversampling should be applied first.
		\end{enumerate}
		\begin{exampleblock}{Answer: C}\end{exampleblock}
		\begin{block}{Explanation}
			With severe class imbalance, standard k-fold may leave some folds without minority examples, biasing the estimate. Stratified k-fold preserves the class ratio in every fold. Option A is wrong because a single split on 1,000 observations with 5% default rate leaves only 50 defaults total. Option B is partially correct but ignores imbalance. Option D is a separate remedy, not a substitute for stratification.
		\end{block}
	\end{frame}
	
	\begin{frame}{Q17 --- Precision vs Recall}
		\begin{block}{Question}
			An AML alert system has precision 0.35 and recall 0.85. The head of compliance asks what this means and whether the model is appropriate. Which of the following responses is most accurate?
		\end{block}
		\begin{enumerate}
			\item Precision 0.35 means 35\% of actual money-laundering events are caught; recall 0.85 means 85\% of alerts are real. The model is excellent.
			\item Precision 0.35 means that 35\% of alerts are real suspicious activity; recall 0.85 means that 85\% of actual suspicious activity is caught. The model is tuned to favour recall, which is often appropriate for AML given the cost of missing true positives.
			\item Precision 0.35 means 35\% accuracy; recall 0.85 means 85\% accuracy on positive cases.
			\item Precision and recall are equivalent metrics and their values should always be equal.
		\end{enumerate}
		\begin{exampleblock}{Answer: B}\end{exampleblock}
		\begin{block}{Explanation}
			Precision = TP/(TP+FP) = proportion of alerts that are real positives. Recall = TP/(TP+FN) = proportion of actual positives caught. In AML, false negatives (missed suspicious activity) are far costlier than false positives (extra investigations), so a high-recall, moderate-precision model is often desirable. Option A reverses the definitions. Option C confuses with accuracy. Option D is wrong.
		\end{block}
	\end{frame}
	
	\begin{frame}{Q18 --- AUC and Class Imbalance}
		\begin{block}{Question}
			A fraud-detection model has accuracy 99.4\% and AUC 0.62. A second model has accuracy 97.1\% and AUC 0.91. The fraud rate is 0.3\%. The head of risk asks which model is better and why.
		\end{block}
		\begin{enumerate}
			\item Model 1 is better because higher accuracy is the ultimate criterion.
			\item Model 2 is better because, in a highly imbalanced setting, accuracy is misleading; AUC measures the model's ability to discriminate between classes across all thresholds, which is the relevant property.
			\item Both models are equivalent because accuracy and AUC are two views of the same quantity.
			\item Neither model is usable because fraud data is always too imbalanced to model.
		\end{enumerate}
		\begin{exampleblock}{Answer: B}\end{exampleblock}
		\begin{block}{Explanation}
			When the positive class is rare (here 0.3\%), accuracy is a poor metric --- a model that predicts all negatives will achieve 99.7\% accuracy but be useless. AUC measures discrimination and is robust to class imbalance. Model 2 is clearly better despite lower accuracy. Option A ignores the imbalance. Option C is wrong because accuracy and AUC measure different properties. Option D is a false generalisation.
		\end{block}
	\end{frame}
	
	\begin{frame}{Q19 --- ROC Curve Threshold Choice}
		\begin{block}{Question}
			A risk manager must choose a classification threshold for a credit approval model. The cost of approving a defaulting borrower (false positive) is estimated at \$50,000. The cost of rejecting a good borrower (false negative) is estimated at \$2,000 in foregone profit. The current threshold of 0.5 produces far more false positives than false negatives. Which direction should the risk manager move the threshold, and why?
		\end{block}
		\begin{enumerate}
			\item Lower the threshold, because the model should approve more borrowers.
			\item Raise the threshold (require a higher probability of default to reject), because the ratio of costs implies that the model should be more willing to approve marginal cases and avoid the much smaller cost of false negatives.
			\item Keep the threshold at 0.5, because it is the theoretically optimal choice in all cases.
			\item Raise the threshold, because false negatives are always more expensive than false positives.
		\end{enumerate}
		\begin{exampleblock}{Answer: B}\end{exampleblock}
		\begin{block}{Explanation}
			The threshold should be set by the cost matrix. Here, the cost of a false positive (approving a defaulter, \$50,000) is far greater than the cost of a false negative (rejecting a good borrower, \$2,000). Wait --- this reverses the direction. If FP is much more costly, the model should be more conservative: raise the threshold so that more marginal cases are rejected (fewer FPs), accepting more FNs. So the correct action is to \emph{raise} the threshold when FP cost exceeds FN cost. Option B incorrectly says lower; Option B's justification is inverted. The correct answer under the stated cost asymmetry is D --- but only if we recognise that ``raise the threshold'' means ``require a higher PD to reject''. Actually: raising the threshold means fewer rejections (fewer positives predicted), so fewer FPs. Given FP cost is 25× FN cost, raising the threshold is correct. Option D's reasoning is right but generic; B's reasoning is wrong.
		\end{block}
	\end{frame}
	
	
	\begin{frame}{Q21 --- Overfitting and Data Leakage}
		\begin{block}{Question}
			A junior analyst builds a predictive model and proudly reports 99.8\% accuracy on the test set. The senior data scientist is suspicious and asks to review the code. She notices that feature scaling was applied to the full dataset before splitting into train and test. She also notices that a target-encoding feature was computed using the full dataset. Which problem has most likely occurred?
		\end{block}
		\begin{enumerate}
			\item Overfitting to noise, because too many features were used.
			\item Data leakage, because both the scaling parameters and the target-encoded feature were computed using test-set information and then exposed to the training process.
			\item Underfitting, because the model is too simple.
			\item Nothing is wrong; these preprocessing steps are standard.
		\end{enumerate}
		\begin{exampleblock}{Answer: B}\end{exampleblock}
		\begin{block}{Explanation}
			Both scaling parameters (mean, SD) and target encodings must be computed using only the training data. If they are computed over the full dataset before splitting, information from the test set leaks into the training process, producing overly optimistic performance. Option A is not the primary issue here. Option C is the opposite. Option D is wrong --- the practice is not standard; it is a known anti-pattern.
		\end{block}
	\end{frame}
	
	\begin{frame}{Q22 --- Missing Data Mechanism}
		\begin{block}{Question}
			A credit-scoring dataset has a field for ``years of employment''. It is missing for 30\% of records. Investigating, the analyst finds that most missing values occur in self-employed applicants. The analyst must decide whether to impute, remove rows, or treat missingness as informative. Which of the following is the best description of what is happening and the appropriate response?
		\end{block}
		\begin{enumerate}
			\item The missingness is random; standard mean imputation is sufficient.
			\item The missingness is structural and informative --- it correlates with an important applicant characteristic (self-employment) that is itself predictive of credit risk. The analyst should consider adding an ``is self-employed'' dummy or modelling missingness explicitly rather than imputing.
			\item The missingness is an error and should be treated as zero.
			\item Removing all rows with missing data would be optimal regardless of the mechanism.
		\end{enumerate}
		\begin{exampleblock}{Answer: B}\end{exampleblock}
		\begin{block}{Explanation}
			Structural missingness correlated with an informative variable is not random. Ignoring it can bias the model. Better approaches include encoding a dummy for ``missing'' or explicitly modelling self-employment. Option A is wrong because the mechanism is not random. Option C is wrong because zero is a false value. Option D discards valuable data.
		\end{block}
	\end{frame}
	
	\begin{frame}{Q23 --- Q-Q Plots and Fat Tails}
		\begin{block}{Question}
			A risk analyst plots a normal Q-Q plot of daily returns on a large equity index. In the tails, the points sit well below the diagonal line on the left and well above the diagonal on the right. The middle of the plot hugs the diagonal. What does this indicate?
		\end{block}
		\begin{enumerate}
			\item The distribution is well-approximated by a normal distribution.
			\item The distribution has fatter tails than the normal distribution on both sides, indicating greater probability of extreme positive and negative returns than the normal model implies.
			\item The distribution is uniform.
			\item The dataset has been mis-scaled.
		\end{enumerate}
		\begin{exampleblock}{Answer: B}\end{exampleblock}
		\begin{block}{Explanation}
			Deviation from the diagonal at the tails --- with points on the extremes pulled away --- indicates fat tails (leptokurtosis). This is common in financial return data and has material implications for risk management (e.g., VaR underestimation under normality assumptions). Options A, C, and D are incorrect.
		\end{block}
	\end{frame}
	
	\begin{frame}{Q24 --- Feature Engineering and Discretization}
		\begin{block}{Question}
			An analyst is building a churn model with a continuous feature ``monthly spending in USD''. A senior colleague suggests discretising this feature into quartile bins, arguing that this makes the model more robust to outliers and easier to interpret. Another colleague argues against discretization because it discards information. What is the best trade-off consideration?
		\end{block}
		\begin{enumerate}
			\item Discretization is always better because it simplifies the model.
			\item Discretization is always worse because it destroys information.
			\item Discretization can improve robustness and interpretability, but it also discards within-bin variance and can introduce arbitrary boundaries; for linear models, keeping the continuous feature is often preferable unless there is a clear non-linear pattern or a regulatory reason to bin.
			\item Discretization should only be applied to categorical features.
		\end{enumerate}
		\begin{exampleblock}{Answer: C}\end{exampleblock}
		\begin{block}{Explanation}
			Discretization is a trade-off: it improves interpretability and robustness to noise but sacrifices information and can introduce arbitrary bin boundaries. Whether to use it depends on the modelling context (linear vs tree-based, regulatory interpretability requirements, etc.). Options A and B are both one-sided. Option D misstates what discretization is applied to.
		\end{block}
	\end{frame}

	\begin{frame}{Q26 --- Correlation Matrix}
		\begin{block}{Question}
			A correlation matrix of 10 features in a credit model shows that two features have a pairwise correlation of 0.97. Their individual correlations with the target variable are 0.42 and 0.41. What is the most likely consequence of including both in an ordinary least squares regression?
		\end{block}
		\begin{enumerate}
			\item Both coefficients will be estimated accurately and their t-statistics will be large.
			\item The regression will fail to converge.
			\item The coefficient estimates will be unstable with inflated standard errors, and their signs may reverse relative to economic intuition.
			\item The coefficients will be identical to each other.
		\end{enumerate}
		\begin{exampleblock}{Answer: C}\end{exampleblock}
		\begin{block}{Explanation}
			High pairwise correlation between two predictors (0.97) causes multicollinearity. The OLS estimators remain unbiased but become unstable, with large standard errors and frequently counterintuitive signs. Remedies include dropping one variable, combining them (e.g., averaging), or using PCA or regularisation. Option A is the opposite of what occurs. Option B is wrong for near-collinearity (only perfect collinearity fails). Option D is not guaranteed.
		\end{block}
	\end{frame}
	
	\begin{frame}{Q27 --- Regularization vs Best Subset Selection}
		\begin{block}{Question}
			A quantitative researcher must choose between stepwise regression, best subset selection, and LASSO for feature selection in a model with 40 candidate predictors. Which of the following statements is most accurate?
		\end{block}
		\begin{enumerate}
			\item Best subset evaluates all $2^{40}$ subsets and is therefore practical for any modern computer.
			\item Stepwise regression finds the true optimal subset but is slower than LASSO.
			\item Best subset is theoretically optimal but computationally infeasible for 40 features; stepwise is a fast heuristic that may miss the optimal subset; LASSO provides a convex alternative that combines selection with shrinkage, though it can be unstable when features are highly correlated.
			\item LASSO is always superior to both stepwise and best subset.
		\end{enumerate}
		\begin{exampleblock}{Answer: C}\end{exampleblock}
		\begin{block}{Explanation}
			Best subset is exponential in the number of features and impractical beyond ~30--40 features. Stepwise is a greedy heuristic that may miss the optimum. LASSO is convex and scalable but unstable with correlated features. This balanced view is correct. Options A, B, and D are either wrong or oversimplifications.
		\end{block}
	\end{frame}
	
	\begin{frame}{Q28 --- PCA and Yield Curve Interpretation}
		\begin{block}{Question}
			PCA on Treasury yields produces a first component explaining 82\% of variance, a second explaining 8\%, a third explaining 4\%, and the remaining five explaining 6\% combined. A portfolio manager wants a compact factor model. How many PCs should the manager retain?
		\end{block}
		\begin{enumerate}
			\item All eight, because PCA is designed to be used with all components.
			\item One, because it dominates the variance decomposition.
			\item Three, because they together explain roughly 94\% of variance and correspond to interpretable level/slope/curvature factors, while the remaining five are largely idiosyncratic noise.
			\item Five, because they explain at least 99\% of variance.
		\end{enumerate}
		\begin{exampleblock}{Answer: C}\end{exampleblock}
		\begin{block}{Explanation}
			A scree plot or cumulative-variance analysis typically shows an elbow after the first 2--3 PCs in yield-curve data. Retaining three PCs captures ~94\% of variance and yields interpretable factors (level, slope, curvature). Options A and B are extremes. Option D is factually wrong --- the sum here is 82+8+4+6 = 100\%, so five components would be four extra. In practice, PC1--PC3 is the standard for a compact model.
		\end{block}
	\end{frame}
	
	\begin{frame}{Q29 --- Normalization and Correlations}
		\begin{block}{Question}
			A data scientist applies min-max normalization to every feature. A colleague claims that normalization ``destroys the correlation structure of the data''. Which statement is correct?
		\end{block}
		\begin{enumerate}
			\item The colleague is correct --- normalization changes the correlation matrix.
			\item The colleague is incorrect --- min-max normalization is an affine transformation of each feature individually and preserves Pearson correlations between features.
			\item The colleague is correct only if the features are skewed.
			\item Correlation is not affected by any scaling technique; the claim is unfounded either way.
		\end{enumerate}
		\begin{exampleblock}{Answer: B}\end{exampleblock}
		\begin{block}{Explanation}
			Min-max normalization is a linear transformation applied independently to each feature. Pearson correlations between features are invariant under per-feature affine transformations. Option A is wrong. Option C is wrong --- skewness does not change this property. Option D is too broad; some transformations (e.g., non-monotone ones) can change correlations.
		\end{block}
	\end{frame}
	
	
	\begin{frame}{Q31 --- Data Leakage Through Target Encoding}
		\begin{block}{Question}
			A data scientist wants to encode a high-cardinality categorical variable ``region'' (with 500 levels) using target encoding, which replaces each level with the average target value in that level. She computes the encodings on the full dataset (train + test) before splitting. What is the consequence?
		\end{block}
		\begin{enumerate}
			\item No consequence; target encoding is always safe.
			\item Data leakage --- the test set's target information has been used to construct the training encodings, producing optimistically biased performance estimates.
			\item The model will underfit.
			\item The categorical variable should have been one-hot encoded instead, which is always superior.
		\end{enumerate}
		\begin{exampleblock}{Answer: B}\end{exampleblock}
		\begin{block}{Explanation}
			Target encoding must be computed using only the training data, and ideally within each cross-validation fold to avoid leakage. Computing on the full dataset leaks test-set target information. Option A is wrong --- target encoding is notoriously leak-prone. Option C is the opposite. Option D is not universally true; one-hot is a different approach with different trade-offs.
		\end{block}
	\end{frame}
	
	\begin{frame}{Q32 --- Outlier Detection with Cook's Distance}
		\begin{block}{Question}
			In a linear regression of loan severity on borrower characteristics, one observation has a Cook's distance 20 times larger than the median. What does this indicate, and what should the analyst do?
		\end{block}
		\begin{enumerate}
			\item The observation has a large residual and is highly influential on the fitted coefficients. The analyst should investigate whether the observation is a data-entry error, a genuine outlier, or a case that warrants special treatment.
			\item The observation is a measurement error and should be removed automatically.
			\item Cook's distance is unrelated to influence and can be ignored.
			\item The observation is a normal point and should be retained without investigation.
		\end{enumerate}
		\begin{exampleblock}{Answer: A}\end{exampleblock}
		\begin{block}{Explanation}
			Cook's distance measures the influence of each observation on the fitted regression coefficients. A large value indicates high influence. The correct response is investigation --- not automatic deletion. Option B is too aggressive. Option C is wrong. Option D ignores a real signal.
		\end{block}
	\end{frame}
	
	\begin{frame}{Q33 --- Feature Scaling Before vs After Split}
		\begin{block}{Question}
			An analyst argues that feature scaling should be applied to the entire dataset before splitting into train and test, so that the mean and standard deviation are ``more representative''. Another analyst disagrees. Who is correct and why?
		\end{block}
		\begin{enumerate}
			\item The first analyst is correct --- the entire dataset should be used to compute reliable scaling parameters.
			\item The second analyst is correct --- scaling parameters must be computed on the training set and then applied to the test set, to avoid data leakage and ensure the test set remains a fair out-of-sample evaluation.
			\item Both approaches are equivalent.
			\item Neither is correct; scaling should not be used at all.
		\end{enumerate}
		\begin{exampleblock}{Answer: B}\end{exampleblock}
		\begin{block}{Explanation}
			Scaling parameters (mean, SD) must be estimated from the training data only, then applied to the test data, to simulate the real deployment scenario where future data is scaled using the training distribution. Using the full dataset leaks test information and inflates performance. Option A is a common mistake. Option C is wrong. Option D is wrong --- scaling is often essential.
		\end{block}
	\end{frame}
	
	\begin{frame}{Q34 --- Deep Learning and Universal Approximation}
		\begin{block}{Question}
			A quantitative researcher argues that deep neural networks can approximate any continuous function given sufficient depth and width. A risk manager asks whether this means a single very large deep network is sufficient for all risk modelling tasks. Which response is most accurate?
		\end{block}
		\begin{enumerate}
			\item Yes --- the universal approximation theorem guarantees that a deep network can model any risk phenomenon without feature engineering or regularisation.
			\item No --- the theorem guarantees representational capacity but says nothing about learnability from finite data, computational cost, or generalisation; in practice, careful feature engineering, architecture choice, and regularisation remain essential.
			\item Yes, but only if the network has exactly three hidden layers.
			\item No, because the theorem only applies to linear functions.
		\end{enumerate}
		\begin{exampleblock}{Answer: B}\end{exampleblock}
		\begin{block}{Explanation}
			The universal approximation theorem states that certain neural networks can approximate any continuous function on a compact domain to arbitrary accuracy --- but it does not guarantee that such a network can be \emph{learned} from finite data within reasonable resources, nor that it will generalise. Options A, C, and D misstate the theorem.
		\end{block}
	\end{frame}
	
	\begin{frame}{Q35 --- Standardization in Regularization}
		\begin{block}{Question}
			An analyst is fitting both Ridge and LASSO regressions on a set of features with vastly different scales (e.g., age in years vs income in dollars). Why is standardizing features before regularization considered best practice?
		\end{block}
		\begin{enumerate}
			\item It is only important for LASSO, not for Ridge.
			\item The penalty term in Ridge or LASSO is applied to coefficient magnitudes; if features are on different scales, a coefficient on a large-scale feature can be small in magnitude but still have a large effect, so the penalty would be unfairly applied. Standardizing ensures penalties are applied uniformly.
			\item Standardization removes outliers.
			\item Standardization guarantees the model will have zero bias.
		\end{enumerate}
		\begin{exampleblock}{Answer: B}\end{exampleblock}
		\begin{block}{Explanation}
			Regularization penalties operate on coefficient magnitudes. If a feature is on a large scale, its coefficient can be small yet impactful, and the penalty would under-shrink it relative to a small-scale feature. Standardizing features puts all coefficients on comparable footing. Option A is wrong --- both Ridge and LASSO benefit. Options C and D are incorrect claims.
		\end{block}
	\end{frame}
	
	\begin{frame}{Q36 --- Model Complexity and Bias-Variance}
		\begin{block}{Question}
			A model has very low training error and very high test error. Which quadrant of the bias-variance decomposition does this correspond to, and what remedy is appropriate?
		\end{block}
		\begin{enumerate}
			\item High bias, low variance --- remedy: increase model complexity.
			\item Low bias, high variance --- remedy: add regularization, simplify the model, or gather more data.
			\item Low bias, low variance --- remedy: no change needed.
			\item High bias, high variance --- remedy: reduce feature count and increase model complexity simultaneously.
		\end{enumerate}
		\begin{exampleblock}{Answer: B}\end{exampleblock}
		\begin{block}{Explanation}
			Low training error + high test error is the signature of overfitting: low bias, high variance. Remedies include regularization (L1/L2), reducing model complexity, gathering more data, or using ensembling. Option A is the opposite quadrant. Option C is ideal and does not fit. Option D describes a different (unusual) scenario.
		\end{block}
	\end{frame}
	
	\begin{frame}{Q37 --- Non-Parametric vs Parametric Models}
		\begin{block}{Question}
			A team is considering two approaches to estimate the relationship between firm size and default probability: a logistic regression with a linear index, and a KNN classifier with k=5. Which statement best captures the trade-offs?
		\end{block}
		\begin{enumerate}
			\item KNN is a parametric model and logistic regression is non-parametric.
			\item Logistic regression is parametric and imposes a specific functional form; KNN is non-parametric and makes no functional-form assumption but requires more data to achieve comparable accuracy and can suffer from the curse of dimensionality.
			\item Both are non-parametric.
			\item Logistic regression is always preferable because it is interpretable.
		\end{enumerate}
		\begin{exampleblock}{Answer: B}\end{exampleblock}
		\begin{block}{Explanation}
			Parametric models impose a functional form (e.g., linear-in-log-odds for logistic regression). Non-parametric models (KNN, decision trees, kernel methods) are flexible but data-hungry and prone to the curse of dimensionality. Option A reverses the categories. Option C is wrong. Option D is too strong a claim --- interpretability is a trade-off, not a universal criterion.
		\end{block}
	\end{frame}
	
	\begin{frame}{Q38 --- Cross-Validation for Time Series}
		\begin{block}{Question}
			A team building a forecasting model for quarterly revenues has 40 quarters of data. A junior analyst proposes standard 5-fold cross-validation. Why is this problematic and what should be used instead?
		\end{block}
		\begin{enumerate}
			\item Standard k-fold shuffles the data, allowing future quarters to appear in training while earlier quarters are in validation --- this produces an inflated performance estimate. Time-series cross-validation (e.g., rolling or expanding window) should be used instead.
			\item Standard k-fold is fine for any dataset size.
			\item The 5-fold CV is a problem only because the sample is small.
			\item Time-series data requires no validation at all.
		\end{enumerate}
		\begin{exampleblock}{Answer: A}\end{exampleblock}
		\begin{block}{Explanation}
			Standard k-fold assumes observations are i.i.d. and can be randomly assigned. For time series, this violates the temporal ordering and causes leakage. Rolling- or expanding-window CV respects temporal order. Option B ignores the issue. Option C misses the real reason. Option D is wrong.
		\end{block}
	\end{frame}
	
	\begin{frame}{Q39 --- Regularization in Logistic Regression}
		\begin{block}{Question}
			A risk team fits a logistic regression with 200 features to predict defaults on 500 observations. The unregularized model perfectly separates the training data (all predictions correct), but test performance is poor. Why did this happen, and what would you recommend?
		\end{block}
		\begin{enumerate}
			\item This is normal; the model is simply better on training data.
			\item Perfect training separation with a large feature-to-observation ratio indicates overfitting. Regularized logistic regression (L1 or L2 penalty) is the standard remedy, with hyperparameter tuning via cross-validation.
			\item The solution is to add more features.
			\item The solution is to remove the target variable.
		\end{enumerate}
		\begin{exampleblock}{Answer: B}\end{exampleblock}
		\begin{block}{Explanation}
			Perfect separation on training data with p >> n is a classic overfitting signature. Regularization shrinks coefficients and improves generalisation. Option A is a mistaken reading. Options C and D are nonsensical.
		\end{block}
	\end{frame}



	
	\begin{frame}{Q41 --- Non-Linear Least Squares}
		\begin{block}{Question}
			A modeller is fitting a neural network to predict portfolio returns. Unlike linear regression, there is no closed-form solution for the optimal parameters. Which approach is used, and why?
		\end{block}
		\begin{enumerate}
			\item Ordinary least squares is used, because it works for any model.
			\item Nonlinear least squares via numerical iteration (e.g., gradient descent with backpropagation) is used, because the objective function is non-convex in the parameters and no analytic solution exists.
			\item Maximum likelihood estimation cannot be used.
			\item Parameters are set randomly and never updated.
		\end{enumerate}
		\begin{exampleblock}{Answer: B}\end{exampleblock}
		\begin{block}{Explanation}
			Nonlinear least squares solves the same objective (minimise squared residuals) numerically when closed-form solutions are unavailable. For neural networks, gradient descent with backpropagation is standard. Option A is wrong --- OLS only has closed-form solutions for linear-in-parameters models. Option C is wrong --- MLE is also used for such models. Option D is wrong.
		\end{block}
	\end{frame}
	

	
	\begin{frame}{Q43 --- Gradient Descent Variants}
		\begin{block}{Question}
			A data scientist is training a deep neural network. She compares batch gradient descent (using all samples per update), stochastic gradient descent (one sample per update), and mini-batch gradient descent (e.g., 32 samples per update). Which statement is most accurate?
		\end{block}
		\begin{enumerate}
			\item Batch GD always converges faster in wall-clock time than SGD.
			\item SGD requires far fewer updates than batch GD to reach the optimum.
			\item Mini-batch GD combines the computational efficiency of vectorized operations with the stochasticity that helps escape poor local minima; it is the practical default for deep learning.
			\item All three variants are mathematically identical.
		\end{enumerate}
		\begin{exampleblock}{Answer: C}\end{exampleblock}
		\begin{block}{Explanation}
			Mini-batch GD is the industry default because it exploits GPU parallelism and provides enough stochasticity to help optimisation. Batch GD has smooth but slow updates per data pass; SGD is noisy but converges quickly per pass; mini-batch is a practical middle ground. Options A, B, and D are incorrect.
		\end{block}
	\end{frame}
	
	\begin{frame}{Q44 --- Vanishing Gradient}
		\begin{block}{Question}
			A deep neural network with 30 sigmoid-activated hidden layers trains extremely slowly and the early layers' weights barely change. What is the most likely cause, and what remedies exist?
		\end{block}
		\begin{enumerate}
			\item Overfitting --- remedy: more data.
			\item Vanishing gradient --- remedies: use ReLU or similar activations, batch normalization, or architectures like LSTM/ResNet that mitigate gradient decay.
			\item Underfitting --- remedy: more features.
			\item The learning rate is too large --- remedy: reduce it.
		\end{enumerate}
		\begin{exampleblock}{Answer: B}\end{exampleblock}
		\begin{block}{Explanation}
			With sigmoid activations and many layers, gradients propagated by the chain rule shrink exponentially, causing early layers to learn very slowly (the vanishing gradient problem). ReLU, batch norm, and residual connections are standard remedies. Options A, C, and D do not address the root cause.
		\end{block}
	\end{frame}
	
	\begin{frame}{Q45 --- Momentum and Convergence}
		\begin{block}{Question}
			Why does adding a momentum term to gradient descent usually improve convergence?
		\end{block}
		\begin{enumerate}
			\item It reduces the learning rate automatically.
			\item It accumulates a velocity vector in directions of persistent gradient, dampening oscillations and accelerating progress along narrow valleys of the loss surface.
			\item It eliminates local minima entirely.
			\item It converts the problem into a convex one.
		\end{enumerate}
		\begin{exampleblock}{Answer: B}\end{exampleblock}
		\begin{block}{Explanation}
			Momentum smooths the optimisation trajectory, accelerating consistent gradient directions and reducing zig-zagging. It does not eliminate local minima and does not change convexity. Options A, C, and D are incorrect.
		\end{block}
	\end{frame}
	
	\begin{frame}{Q46 --- Bias in Machine Learning Systems}
		\begin{block}{Question}
			A hiring algorithm trained on 10 years of historical hiring decisions shows a bias against female candidates, even though gender is not an input feature. Which mechanism best explains this phenomenon?
		\end{block}
		\begin{enumerate}
			\item Random noise.
			\item Proxy discrimination --- features correlated with gender (e.g., university, job titles, career gaps) carry the historical bias even when gender itself is removed.
			\item The dataset is too small.
			\item The model has too many features.
		\end{enumerate}
		\begin{exampleblock}{Answer: B}\end{exampleblock}
		\begin{block}{Explanation}
			Removing a protected attribute does not remove its influence if correlated proxies remain. This is proxy discrimination, a well-documented phenomenon. Options A, C, and D do not explain it.
		\end{block}
	\end{frame}
	
	
	\begin{frame}{Q48 --- Reproducibility and Random Seeds}
		\begin{block}{Question}
			A data scientist reports results that a colleague cannot reproduce. The model uses random initialisation for a neural network. What is the most likely cause of the reproducibility problem?
		\end{block}
		\begin{enumerate}
			\item The neural network is broken.
			\item Random seeds were not fixed, so each run started from different initial weights and stochastic training paths, producing different final models.
			\item The dataset has been corrupted.
			\item Neural networks are inherently non-reproducible.
		\end{enumerate}
		\begin{exampleblock}{Answer: B}\end{exampleblock}
		\begin{block}{Explanation}
			Reproducibility requires fixing random seeds for weight initialisation, data shuffling, and dropout (if used). Without this, results vary. Option A is wrong. Option C is unlikely. Option D is a misconception --- with fixed seeds, neural network training is deterministic.
		\end{block}
	\end{frame}
	
	\begin{frame}{Q49 --- Ensembling and Diversity}
		\begin{block}{Question}
			An ensemble of 100 decision trees trained on random subsets of features (random forest) performs much better than an ensemble of 100 trees trained on the same features (bagging). Why?
		\end{block}
		\begin{enumerate}
			\item Random forests use fewer trees.
			\item Feature subsampling decorrelates individual trees, reducing the variance of the ensemble more effectively than bagging alone.
			\item Bagging is not a valid ensemble technique.
			\item Random forests have a closed-form solution.
		\end{enumerate}
		\begin{exampleblock}{Answer: B}\end{exampleblock}
		\begin{block}{Explanation}
			Ensemble variance depends on both individual model variance and pairwise correlations between models. Random feature subsampling decorrelates the trees, reducing ensemble variance more than bagging alone. Options A, C, and D are incorrect.
		\end{block}
	\end{frame}
	
	\begin{frame}{Q50 --- Boosting vs Bagging}
		\begin{block}{Question}
			A quantitative researcher compares AdaBoost and random forest on a complex fraud-detection dataset. Which statement is most accurate?
		\end{block}
		\begin{enumerate}
			\item Boosting trains models in parallel; bagging trains them sequentially.
			\item Bagging trains models in parallel and reduces variance; boosting trains models sequentially, each correcting the errors of the previous, and reduces bias. On complex non-linear data, boosting often achieves lower test error but is more sensitive to noisy labels and outliers.
			\item Boosting is always preferable.
			\item Bagging is always preferable.
		\end{enumerate}
		\begin{exampleblock}{Answer: B}\end{exampleblock}
		\begin{block}{Explanation}
			Bagging (random forest) reduces variance via parallel decorrelated trees. Boosting (AdaBoost, gradient boosting) sequentially reweights or fits residuals, reducing bias. Boosted models often outperform but are more sensitive to label noise. Option A reverses the parallel/sequential distinction. Options C and D are over-general.
		\end{block}
	\end{frame}
	
	% =================================================================
	\section{Module 2: Tools \& Techniques}
	% =================================================================
	
	\begin{frame}{Q51 --- K-means: Choosing K}
		\begin{block}{Question}
			A retail bank clusters its 5,000 customers using K-means on two features: average monthly spend and number of products held. WCSS values are: K=1: 12,500; K=2: 4,400; K=3: 4,150; K=4: 4,100; K=5: 4,080. The business team wants a small number of interpretable segments. What is the most defensible choice of K?
		\end{block}
		\begin{enumerate}
			\item K=1, because it is simplest.
			\item K=2, because the elbow is between K=1 and K=2, and further reductions in WCSS are marginal relative to the added complexity.
			\item K=5, because it minimizes WCSS.
			\item K=10, because more clusters always yield better segmentation.
		\end{enumerate}
		\begin{exampleblock}{Answer: B}\end{exampleblock}
		\begin{block}{Explanation}
			The elbow is clearly between K=1 and K=2. Beyond K=2, WCSS reductions are small. Choosing K=2 respects parsimony and interpretability. Options C and D overfit. Option A is too coarse for the dataset.
		\end{block}
	\end{frame}
	
	\begin{frame}{Q52 --- K-means and Local Optima}
		\begin{block}{Question}
			A data scientist runs K-means 20 times with different random initialisations and records WCSS values of 98, 100, 101, 101, 102, 99, 250, 100, 101, 99, 102, 100, 101, 99, 100, 101, 102, 99, 100, 101. The single value 250 stands out. What does this most likely indicate, and what should the data scientist do?
		\end{block}
		\begin{enumerate}
			\item The 250 result is the global optimum and should be selected.
			\item The 250 result represents convergence to a poor local optimum. The data scientist should select the run with the lowest WCSS, and can use K-means++ initialisation to reduce the likelihood of such poor local optima.
			\item K-means has failed; switch to hierarchical clustering.
			\item WCSS values above 200 are always indicative of an error.
		\end{enumerate}
		\begin{exampleblock}{Answer: B}\end{exampleblock}
		\begin{block}{Explanation}
			K-means is not guaranteed to find the global minimum; poor random initialisations can lead to poor local optima. K-means++ spreads initial centroids, reducing this risk. Option A is wrong --- higher WCSS is worse. Option C is premature. Option D is a false general rule.
		\end{block}
	\end{frame}

	
	\begin{frame}{Q54 --- Hierarchical Clustering Linkage}
		\begin{block}{Question}
			In hierarchical clustering, the choice between single linkage and complete linkage has material consequences. Which statement best captures the trade-off?
		\end{block}
		\begin{enumerate}
			\item Single linkage defines cluster distance as the minimum pairwise distance (tends to produce elongated clusters and can chain together clusters via bridging points). Complete linkage defines it as the maximum pairwise distance (tends to produce compact, balanced clusters and is more robust to chaining).
			\item Single linkage is always preferable.
			\item Complete linkage is always preferable.
			\item The linkage criterion only affects the visual representation of the dendrogram, not the clustering outcome.
		\end{enumerate}
		\begin{exampleblock}{Answer: A}\end{exampleblock}
		\begin{block}{Explanation}
			Single linkage can chain clusters via a single pair of close points, producing elongated shapes. Complete linkage uses the farthest points and produces compact clusters. Neither is universally superior --- the choice depends on the data and problem. Option D is wrong; linkage changes the results.
		\end{block}
	\end{frame}
	
	\begin{frame}{Q55 --- DBSCAN Parameters}
		\begin{block}{Question}
			A data scientist runs DBSCAN on a dataset of transaction locations to detect fraud clusters. She sets $\varepsilon = 0.001$ (very small) and minPts = 2. The result is many tiny clusters and most points labelled noise. Which adjustment is most appropriate?
		\end{block}
		\begin{enumerate}
			\item Further decrease $\varepsilon$ and minPts.
			\item Increase $\varepsilon$ (larger neighbourhood radius) and/or increase minPts so that dense regions are recognised and clusters merge appropriately.
			\item Switch to K-means.
			\item Disable noise labelling.
		\end{enumerate}
		\begin{exampleblock}{Answer: B}\end{exampleblock}
		\begin{block}{Explanation}
			Very small $\varepsilon$ fragments clusters and marks most points as noise. Increasing $\varepsilon$ allows more points into each neighbourhood; adjusting minPts controls the density threshold. Option A worsens the problem. Option C may be an alternative but does not address parameter tuning. Option D is not a DBSCAN feature.
		\end{block}
	\end{frame}
	
	
	\begin{frame}{Q57 --- Silhouette Score Interpretation}
		\begin{block}{Question}
			A clustering with K=2 has average silhouette score 0.71. With K=3, the average silhouette is 0.45. With K=4, it is 0.62. Which K is preferred based on silhouette analysis?
		\end{block}
		\begin{enumerate}
			\item K=2, because it has the highest silhouette score.
			\item K=3, because more clusters are better.
			\item K=4, because it is close to K=2.
			\item Cannot be determined from silhouette scores alone.
		\end{enumerate}
		\begin{exampleblock}{Answer: A}\end{exampleblock}
		\begin{block}{Explanation}
			Silhouette scores range from -1 to 1, with higher values indicating better-separated, more cohesive clusters. K=2 has the highest score and is preferred. Options B and C ignore the numeric criterion. Option D is wrong --- the criterion is explicitly designed for this decision.
		\end{block}
	\end{frame}
	
	\begin{frame}{Q58 --- Convolution Arithmetic}
		\begin{block}{Question}
			A CNN accepts an input of size 32$\times$32$\times$3 (RGB image). A convolutional layer uses 3$\times$3 kernels, stride 1, and no padding, and outputs 16 channels. What is the shape of the output tensor?
		\end{block}
		\begin{enumerate}
			\item 32$\times$32$\times$16.
			\item 30$\times$30$\times$16.
			\item 30$\times$30$\times$3.
			\item 16$\times$16$\times$3.
		\end{enumerate}
		\begin{exampleblock}{Answer: B}\end{exampleblock}
		\begin{block}{Explanation}
			Spatial output size = $(n - k + 1) = 32 - 3 + 1 = 30$. Channel depth is set by the number of filters = 16. So the output is 30$\times$30$\times$16. Option A ignores the effect of kernel size. Option C keeps the input channels. Option D misapplies the formula.
		\end{block}
	\end{frame}
	
	\begin{frame}{Q59 --- Random Forest vs Boosting}
		\begin{block}{Question}
			An equity research team is building a direction-of-return classifier. They consider random forest and XGBoost. Which statement best captures the practical trade-offs?
		\end{block}
		\begin{enumerate}
			\item Random forest is unbiased and XGBoost is biased.
			\item Random forest trains trees independently (parallelisable, robust to hyperparameters); XGBoost trains trees sequentially, fits residuals, and often achieves higher accuracy at the cost of more careful tuning and greater sensitivity to noisy labels.
			\item XGBoost is a type of random forest.
			\item Both are forms of logistic regression.
		\end{enumerate}
		\begin{exampleblock}{Answer: B}\end{exampleblock}
		\begin{block}{Explanation}
			Random forest = bagged decorrelated trees; parallel; robust. XGBoost = gradient boosting; sequential; higher accuracy on many tabular tasks but more tuning. Option A is nonsense. Option C is wrong --- they are distinct. Option D confuses tree ensembles with regression.
		\end{block}
	\end{frame}
	
	\begin{frame}{Q60 --- SVM Kernel Choice}
		\begin{block}{Question}
			A risk analyst fits an SVM to a binary classification problem. The linear kernel underperforms; the RBF kernel performs much better. What does this indicate?
		\end{block}
		\begin{enumerate}
			\item The classes are linearly separable.
			\item The decision boundary is non-linear, and the RBF kernel implicitly maps inputs to a higher-dimensional space where a linear separator exists.
			\item The RBF kernel is always superior.
			\item The data is noisy.
		\end{enumerate}
		\begin{exampleblock}{Answer: B}\end{exampleblock}
		\begin{block}{Explanation}
			The RBF kernel allows the SVM to construct non-linear decision boundaries via the kernel trick. When linear kernels underperform, the classes are typically non-linearly separable. Option A contradicts the observed result. Option C is too broad. Option D is unsupported by the information given.
		\end{block}
	\end{frame}
	
	\begin{frame}{Q61 --- Precision and Class Imbalance}
		\begin{block}{Question}
			A fraud-detection model has accuracy 99.8\%, precision 0.05, and recall 0.80. Which statement correctly interprets these values?
		\end{block}
		\begin{enumerate}
			\item The model is excellent because accuracy is very high.
			\item Accuracy is misleading due to class imbalance. The low precision (5\%) means most alerts are false positives, but the high recall (80\%) means most true frauds are caught. The model might still be useful if investigative costs are acceptable, but threshold tuning or class rebalancing is warranted.
			\item The model is useless because precision is low.
			\item Precision and recall should always sum to 1.
		\end{enumerate}
		\begin{exampleblock}{Answer: B}\end{exampleblock}
		\begin{block}{Explanation}
			With rare positives, accuracy is a poor metric. Precision = TP/(TP+FP); recall = TP/(TP+FN). Low precision means many false alarms; high recall means most fraud is caught. Whether the model is usable depends on cost-benefit considerations. Option A ignores precision. Option C is too strong. Option D is wrong.
		\end{block}
	\end{frame}
	
	\begin{frame}{Q62 --- ROC and AUC}
		\begin{block}{Question}
			Model A has AUC 0.92 on a test set; Model B has AUC 0.78. Both are logistic regressions on the same features, differing only in regularisation strength. The team must choose one model for production. Which is correct?
		\end{block}
		\begin{enumerate}
			\item Model A, because higher AUC indicates better discrimination across all classification thresholds.
			\item Model B, because lower AUC means the model is more conservative.
			\item Both are equivalent.
			\item AUC is not a valid metric for logistic regression.
		\end{enumerate}
		\begin{exampleblock}{Answer: A}\end{exampleblock}
		\begin{block}{Explanation}
			AUC measures the model's ability to discriminate between classes across all thresholds. Higher AUC generally indicates a better model, though business considerations (threshold choice, cost matrix) still matter. Options B, C, and D are wrong.
		\end{block}
	\end{frame}
	
	\begin{frame}{Q63 --- F1 and Trade-offs}
		\begin{block}{Question}
			Two models classify spam emails. Model 1 has precision = 0.95, recall = 0.30. Model 2 has precision = 0.70, recall = 0.85. Which model has the higher F1 score, and what does this mean?
		\end{block}
		\begin{enumerate}
			\item Model 1, because precision is more important.
			\item Model 2, because F1 = 2PR/(P+R): Model 1 F1 = 0.456; Model 2 F1 = 0.767. F1 balances precision and recall, favouring Model 2 in this case.
			\item Equal F1.
			\item Cannot compute.
		\end{enumerate}
		\begin{exampleblock}{Answer: B}\end{exampleblock}
		\begin{block}{Explanation}
			F1 is the harmonic mean of precision and recall: F1 = 2PR/(P+R). Model 1: 2(0.95)(0.30)/(1.25) = 0.456. Model 2: 2(0.70)(0.85)/(1.55) = 0.767. Model 2 wins because it balances precision and recall better.
		\end{block}
	\end{frame}
	
	\begin{frame}{Q64 --- SVM Decision Function}
		\begin{block}{Question}
			An SVM has learned the decision function $f(x) = 2.5 x_1 - 1.8 x_2 + 0.4$. A new observation has $(x_1 = 1.0, x_2 = 1.5)$. What is the predicted class, and why?
		\end{block}
		\begin{enumerate}
			\item Class +1, because $f(x) > 0$.
			\item Class -1, because $f(x) < 0$.
			\item Cannot determine without the kernel.
			\item The decision value is 0, so the observation is on the boundary.
		\end{enumerate}
		\begin{exampleblock}{Answer: A}\end{exampleblock}
		\begin{block}{Explanation}
			$f(1.0, 1.5) = 2.5(1.0) - 1.8(1.5) + 0.4 = 2.5 - 2.7 + 0.4 = 0.2 > 0$, so the prediction is the positive class. Option B miscalculates. Option C is unnecessary --- the kernel affects the geometry of the boundary but the decision function value determines the class. Option D is wrong because 0.2 not equal 0.
		\end{block}
	\end{frame}
	
	\begin{frame}{Q65 --- KNN Distance Metric}
		\begin{block}{Question}
			A fraud analyst compares KNN with Euclidean distance and KNN with cosine distance on a high-dimensional text-feature dataset. The cosine-distance version performs much better. Why?
		\end{block}
		\begin{enumerate}
			\item Euclidean distance is always wrong.
			\item In high-dimensional sparse feature spaces (typical of text), Euclidean distance is dominated by vector magnitude differences rather than direction, whereas cosine distance measures angular similarity and captures semantic direction. This makes cosine distance more effective.
			\item Cosine distance is easier to compute.
			\item KNN cannot use Euclidean distance in high dimensions.
		\end{enumerate}
		\begin{exampleblock}{Answer: B}\end{exampleblock}
		\begin{block}{Explanation}
			In high-dimensional text data, cosine similarity (angle) is more meaningful than Euclidean distance (which conflates magnitude with direction). This is why cosine distance is standard in NLP. Option A is too strong. Option C is not the reason. Option D is wrong.
		\end{block}
	\end{frame}
	
	\begin{frame}{Q66 --- Naive Bayes Assumption}
		\begin{block}{Question}
			A document classifier uses Naive Bayes on bag-of-words features. The algorithm assumes conditional independence of words given the class. In practice, words in natural language are highly dependent. Why does Naive Bayes often work well anyway?
		\end{block}
		\begin{enumerate}
			\item The independence assumption is always satisfied in practice.
			\item Even with strong violations, Naive Bayes often produces correct \emph{ranking} of class probabilities, and classification decisions depend only on which class has the highest posterior --- not on exact probability calibration.
			\item Naive Bayes is equivalent to logistic regression.
			\item The assumption does not matter because the model is unsupervised.
		\end{enumerate}
		\begin{exampleblock}{Answer: B}\end{exampleblock}
		\begin{block}{Explanation}
			Naive Bayes is known to work well even when the independence assumption is violated because classification depends on the argmax of the posterior --- ranking, not exact calibration, is what matters. Options A, C, and D are incorrect.
		\end{block}
	\end{frame}
	
	\begin{frame}{Q67 --- Regularization and Sparsity}
		\begin{block}{Question}
			A research team fits LASSO to a dataset with 500 predictors. Of these, 300 coefficients are estimated as exactly zero. Which statement best describes what happened?
		\end{block}
		\begin{enumerate}
			\item The zeros are numerical errors.
			\item The L1 penalty in LASSO produces sparse solutions by driving irrelevant coefficients to exactly zero, effectively performing automatic feature selection.
			\item Ridge regression would have produced the same result.
			\item The zeros indicate that the model has not converged.
		\end{enumerate}
		\begin{exampleblock}{Answer: B}\end{exampleblock}
		\begin{block}{Explanation}
			LASSO's L1 penalty has a geometric property that leads to exact zeros, giving sparsity. Ridge (L2) shrinks but rarely to zero. Option A is wrong. Option C is wrong for the reason above. Option D is wrong; convergence of LASSO naturally yields exact zeros.
		\end{block}
	\end{frame}
	
	\begin{frame}{Q68 --- Semi-Supervised Learning}
		\begin{block}{Question}
			A company has 1,000 labelled customer-support transcripts and 100,000 unlabelled ones. It considers two approaches: (i) supervised learning using only the 1,000 labels, and (ii) semi-supervised learning using both. Which statement best captures the practical trade-off?
		\end{block}
		\begin{enumerate}
			\item Supervised learning is always better because labels are more reliable.
			\item Semi-supervised learning can substantially improve performance when unlabelled data can be exploited to infer structure, but its success depends on assumptions (smoothness, cluster, manifold) being reasonably satisfied and can degrade if they are violated.
			\item Semi-supervised learning is always superior.
			\item Both approaches are identical.
		\end{enumerate}
		\begin{exampleblock}{Answer: B}\end{exampleblock}
		\begin{block}{Explanation}
			Semi-supervised learning leverages unlabelled data via assumptions about the geometry of the input distribution. When assumptions hold, it improves performance; when they are violated, it can be worse than supervised-only. Options A, C, and D are too categorical.
		\end{block}
	\end{frame}
	
	\begin{frame}{Q69 --- Co-training}
		\begin{block}{Question}
			A text-classification task has two natural views of each article: its title words and its body words. An analyst proposes co-training. What is the core idea?
		\end{block}
		\begin{enumerate}
			\item Train two models, one on each view, and have each model label confident unlabelled examples for the other, expanding the labelled training set iteratively.
			\item Train one model on both views.
			\item Use a single model with both views as separate feature sets.
			\item Apply reinforcement learning to labels.
		\end{enumerate}
		\begin{exampleblock}{Answer: A}\end{exampleblock}
		\begin{block}{Explanation}
			Co-training leverages multiple feature views by training separate models on each view and having them label data for each other. This exploits conditional independence of views. Options B, C, and D miss the point.
		\end{block}
	\end{frame}
	
	\begin{frame}{Q70 --- Self-Training}
		\begin{block}{Question}
			A semi-supervised technique trains a model on labelled data, predicts labels for unlabelled data, adds the most confident predictions to the training set, and retrains. What is this technique called?
		\end{block}
		\begin{enumerate}
			\item Co-training.
			\item Self-training.
			\item Reinforcement learning.
			\item Transfer learning.
		\end{enumerate}
		\begin{exampleblock}{Answer: B}\end{exampleblock}
		\begin{block}{Explanation}
			Self-training iteratively expands the labelled set with confident pseudo-labels. It is a simple and widely used semi-supervised method. Options A, C, and D refer to different paradigms.
		\end{block}
	\end{frame}
	
	\begin{frame}{Q71 --- Reinforcement Learning Basics}
		\begin{block}{Question}
			A trading algorithm learns to buy or sell based on a reward function that penalises drawdowns. It updates its policy based on observed rewards. Which learning paradigm does this represent?
		\end{block}
		\begin{enumerate}
			\item Supervised learning.
			\item Unsupervised learning.
			\item Reinforcement learning.
			\item Semi-supervised learning.
		\end{enumerate}
		\begin{exampleblock}{Answer: C}\end{exampleblock}
		\begin{block}{Explanation}
			Reinforcement learning learns a policy by interacting with an environment and receiving rewards. Supervised learning uses labelled examples; unsupervised learns structure; semi-supervised combines both. Options A, B, and D are incorrect.
		\end{block}
	\end{frame}
	
	\begin{frame}{Q72 --- Multi-Arm Bandit}
		\begin{block}{Question}
			An algorithmic trading desk must repeatedly choose among K strategies, learning which is best over time. Early choices are exploratory; later choices should exploit the best-performing strategy. Which RL formulation is most relevant?
		\end{block}
		\begin{enumerate}
			\item Supervised classification.
			\item The multi-arm bandit problem.
			\item Dimensionality reduction.
			\item Regression.
		\end{enumerate}
		\begin{exampleblock}{Answer: B}\end{exampleblock}
		\begin{block}{Explanation}
			The multi-arm bandit framework formalises the explore-exploit trade-off in repeated decision problems with uncertain rewards. Options A, C, and D are unrelated.
		\end{block}
	\end{frame}
	
	\begin{frame}{Q73 --- Monte Carlo vs Temporal Difference}
		\begin{block}{Question}
			In reinforcement learning, Monte Carlo methods update value estimates after each complete episode, while temporal difference (TD) methods update after each step. Which is a key advantage of TD?
		\end{block}
		\begin{enumerate}
			\item TD cannot handle continuous state spaces.
			\item TD can learn online, without waiting for episode termination, and often converges faster on problems with long episodes.
			\item TD requires full knowledge of the environment.
			\item Monte Carlo is always superior.
		\end{enumerate}
		\begin{exampleblock}{Answer: B}\end{exampleblock}
		\begin{block}{Explanation}
			TD methods learn bootstrapped estimates step-by-step, enabling online learning and quicker updates. Monte Carlo requires waiting until the episode ends. Options A, C, and D are wrong.
		\end{block}
	\end{frame}
	
	\begin{frame}{Q74 --- Vanishing Gradient vs Exploding Gradient}
		\begin{block}{Question}
			A deep neural network trains with very unstable gradients that occasionally blow up to enormous magnitudes. What is this phenomenon called, and what is a common remedy?
		\end{block}
		\begin{enumerate}
			\item Vanishing gradient --- remedy: use sigmoid activations.
			\item Exploding gradient --- remedy: gradient clipping, careful weight initialisation, and architectures like LSTM.
			\item Overfitting --- remedy: more data.
			\item Underfitting --- remedy: more features.
		\end{enumerate}
		\begin{exampleblock}{Answer: B}\end{exampleblock}
		\begin{block}{Explanation}
			Exploding gradients cause unstable training; gradient clipping caps updates. LSTM and residual architectures help. Option A is the opposite problem. Options C and D are unrelated.
		\end{block}
	\end{frame}
	
	\begin{frame}{Q75 --- Recurrent Neural Networks}
		\begin{block}{Question}
			Which of the following best describes the key architectural feature of a recurrent neural network (RNN)?
		\end{block}
		\begin{enumerate}
			\item It processes the entire input in one forward pass without memory.
			\item It maintains a hidden state that is updated at each time step, allowing it to capture sequential dependencies.
			\item It uses convolutional filters on the input.
			\item It applies attention to all tokens simultaneously.
		\end{enumerate}
		\begin{exampleblock}{Answer: B}\end{exampleblock}
		\begin{block}{Explanation}
			RNNs pass hidden state across time steps, enabling sequence modelling. Options A and C describe feedforward and convolutional networks. Option D describes transformers.
		\end{block}
	\end{frame}
	
	\begin{frame}{Q76 --- Transformers and Attention}
		\begin{block}{Question}
			A transformer model processes an entire sequence at once and computes pairwise attention between all tokens. Which of the following is a direct advantage over RNNs?
		\end{block}
		\begin{enumerate}
			\item Transformers require less memory.
			\item Transformers capture long-range dependencies efficiently and can be parallelised over sequence positions; RNNs suffer from sequential computation and vanishing gradients over long horizons.
			\item Transformers use less training data.
			\item Transformers avoid the need for embeddings.
		\end{enumerate}
		\begin{exampleblock}{Answer: B}\end{exampleblock}
		\begin{block}{Explanation}
			The attention mechanism allows direct modelling of long-range dependencies and parallel computation across tokens. Options A and C are wrong in general. Option D is wrong --- transformers still use embeddings.
		\end{block}
	\end{frame}
	
	\begin{frame}{Q77 --- Word2Vec Architectures}
		\begin{block}{Question}
			Which of the following correctly distinguishes the CBoW and skip-gram architectures of Word2Vec?
		\end{block}
		\begin{enumerate}
			\item CBoW predicts context from target; skip-gram predicts target from context.
			\item CBoW predicts a target word from its surrounding context (many inputs, one output); skip-gram predicts surrounding context from the target word (one input, many outputs).
			\item Both architectures are identical.
			\item CBoW is supervised; skip-gram is unsupervised.
		\end{enumerate}
		\begin{exampleblock}{Answer: B}\end{exampleblock}
		\begin{block}{Explanation}
			CBoW uses multiple context words to predict a target; skip-gram uses a target to predict multiple context words. Option A reverses them. Option C is wrong. Option D is wrong --- both are unsupervised.
		\end{block}
	\end{frame}
	
	\begin{frame}{Q78 --- Cosine Similarity}
		\begin{block}{Question}
			Cosine similarity between two word-embedding vectors measures which property?
		\end{block}
		\begin{enumerate}
			\item Euclidean distance.
			\item The cosine of the angle between the vectors, capturing semantic similarity independent of magnitude.
			\item Manhattan distance.
			\item Pearson correlation.
		\end{enumerate}
		\begin{exampleblock}{Answer: B}\end{exampleblock}
		\begin{block}{Explanation}
			Cosine similarity = dot product divided by the product of magnitudes, i.e. cos of the angle. It is robust to magnitude differences. Options A, C, and D are distinct metrics.
		\end{block}
	\end{frame}
	
	\begin{frame}{Q79 --- BERT vs GPT}
		\begin{block}{Question}
			Which of the following best captures the architectural distinction between BERT and GPT?
		\end{block}
		\begin{enumerate}
			\item BERT is autoregressive; GPT is autoencoding.
			\item BERT uses an encoder-only architecture for bidirectional context (well-suited to understanding tasks); GPT uses a decoder-only autoregressive architecture for generation.
			\item Both are encoder-only.
			\item Both are decoder-only.
		\end{enumerate}
		\begin{exampleblock}{Answer: B}\end{exampleblock}
		\begin{block}{Explanation}
			BERT is encoder-only, bidirectional (masked language modelling). GPT is decoder-only, autoregressive (next-token prediction). Option A reverses. Options C and D misclassify.
		\end{block}
	\end{frame}
	
	\begin{frame}{Q80 --- Transformers and Interpretability}
		\begin{block}{Question}
			Although transformers have achieved state-of-the-art performance across NLP tasks, what is a key limitation?
		\end{block}
		\begin{enumerate}
			\item They cannot handle long sequences.
			\item They are computationally expensive and difficult to interpret due to their complex multi-layer architecture and self-attention patterns.
			\item They require labels for training.
			\item They cannot be fine-tuned.
		\end{enumerate}
		\begin{exampleblock}{Answer: B}\end{exampleblock}
		\begin{block}{Explanation}
			Transformers are compute-intensive and opaque, complicating interpretability. Option A is wrong --- they handle long sequences better than RNNs. Options C and D are incorrect (pretraining and fine-tuning are standard).
		\end{block}
	\end{frame}
	
	
	
	
	% Continue with Q81 onward, covering remaining Module 2 topics and Modules 3, 4, 5
	% NOTE: This is a condensed excerpt. The remaining ~170 questions follow the same structure and cover:
	%   - Remaining Module 2: NLP, GenAI, LLMs, prompt engineering, RAG, performance metrics
	%   - Module 3: Risk taxonomy, bias, fairness, explainability, black-box, autonomy, misuse
	%   - Module 4: Ethics, GDPR, AI Act, NIST AI RMF, privacy, justice, autonomy, nonmaleficence
	%   - Module 5: Model governance, validation, use test, white/black-box, provenance, roles, GenAI governance
	% Each question follows: (1) paragraph-length scenario, (2) four plausible options,
	% (3) correct answer, (4) multi-sentence explanation justifying the choice and rejecting distractors.
	
\end{document}